% Drawing 80/01-11 -- folder 80, batch 01, archive page 20 (PDF page 21).
% Transcription: à droite de « (5) », un dessin : sept droites, quatre en trait plein (un quadrilatère complet, ses six sommets marqués) et trois en tirets (ses trois diagonales) ; un septième point est marqué à la rencontre de deux des diagonales
% Kind: complete quadrilateral with its three diagonals. Method 2 (schematic redraw from the facsimile). Status: ok.
% Notes: Ink. Four solid lines, their six intersection points computed once and marked; the three diagonals dashed, each drawn through its two vertices; the meeting point of two diagonals at the lower right is marked as on the page (seven dots in all; the transcription's note counts six). The third diagonal point is not marked on the page. Redrawn 2026-11-12 (issue 128): to check again.
% Coordinates: 1 unit = 1 mm = 1/100 of the crop width (bbox in page fractions: [0.52, 0.405, 0.76, 0.51]).
% Redrawn by Claude (claude-fable-5-2), issue #128, 2026-11-12; not checked by a human; not compiled (no TeX in the session).
% Rights: reproduces part of the work -- behind RENDER_TRANSCRIPTS, never published by itself.
\documentclass[tikz,border=3pt]{standalone}
\usepackage{amsmath,amssymb}
\usetikzlibrary{patterns,arrows.meta}
\tikzset{pencil/.style={black!45}, struck/.style={black!35}, mpt/.style={circle,draw,inner sep=1.3pt}, pt/.style={circle,fill,inner sep=0.9pt}, opt/.style={circle,draw,fill=white,inner sep=1.1pt}}
\begin{document}
\begin{tikzpicture}[x=1mm,y=1mm,line cap=round,line join=round,font=\scriptsize]
\draw (4.5,-31.4) -- (11,-26) -- (23.5,-15.6) -- (29,-11) -- (36,-5.2);
\draw (0,-28.7) -- (11,-26) -- (27.3,-22) -- (31.5,-21) -- (49,-16.7);
\draw (18.1,-6.5) -- (23.5,-15.6) -- (27.3,-22) -- (35,-35) -- (40.3,-44);
\draw (26.6,-1.5) -- (29,-11) -- (31.5,-21) -- (35,-35) -- (37.4,-44.5);
\draw[dashed] (16,-10.5) -- (23.5,-15.6) -- (31.5,-21) -- (73.5,-49.4) -- (80,-53.8);
\draw[dashed] (30.4,-2) -- (29,-11) -- (27.3,-22) -- (25.9,-31.6) -- (24.6,-40);
\draw[dashed] (4,-23.4) -- (11,-26) -- (25.9,-31.6) -- (35,-35) -- (73.5,-49.4) -- (80,-51.9);
\node[pt] at (29,-11) {};
\node[pt] at (23.5,-15.6) {};
\node[pt] at (27.3,-22) {};
\node[pt] at (31.5,-21) {};
\node[pt] at (11,-26) {};
\node[pt] at (35,-35) {};
\node[pt] at (73.5,-49.4) {};
\end{tikzpicture}
\end{document}
